\section{Conclusion}
\label{sec:conclusion}

This study reports a \DeploymentDays{}-day agricultural edge-IoT deployment and separately executed live experiments examining the historical implementation's performance, all recomputed for this version directly from an analyst-supplied raw-data package rather than transcribed from an earlier draft. An implementation-defined authorization-associated span averaged \SpanLatencyPeerFour{} ms within the \TotalLatencyPeerFour{} ms write path. Across \PeerRunsPerConfig{} counterbalanced runs per configuration, confirmed against the released run manifest, increasing logical peer multiplicity from 4 to 32 on four fixed hosts was associated with a \RemainingLatencyDiffFourVThirtyTwo{} ms reduction in the uninstrumented remainder, while the recorded span differed by only \SpanDiffFourVThirtyTwo{} ms between the extreme configurations. The experiment measured latency rather than physical scale-out or peer-count throughput and did not identify the internal mechanism responsible for the association.

A separate interleaved experiment found that the complete access-control-enabled path averaged \ThroughputPctDiffMean\% lower throughput than the comparison path. No condition-by-concurrency interaction was detected across the tested closed-loop range, but the experiment does not establish a universal or constant authorization penalty. Revocation logs contained \RevocationStageRecords{} stage records, including \RevocationDelayedCount{} carrying a delayed label, but the trace links stages only to a single event per subject, preventing reconstruction of end-to-end revocation exposure or causal attribution.

Code-to-specification audit, reverified directly against the released source for this version, identified three cross-layer defect classes: authorization-policy deviations comprising unintended audit inheritance and Agronomist actuation permission; invalid two-residue reconstruction for \CrtTwoResidueCount{} readings; and signature-verification failure comprising a signed-message mismatch and fail-open acceptance of missing signatures. It also identified a revocation-observability gap. Corrections to each defect are specified and arithmetically checked, but -- unlike an earlier internal draft's description -- no corrected chaincode, firmware or gateway package exists in the repository analyzed here; transfer to any future corrected implementation is therefore not established, and no corrected-version performance or security claim is made. The released raw measurements, execution schedules, the disclosed security oracle, the audit evidence and the analysis pipeline that recomputes every number in this article allow the findings and their limitations to be independently examined.
